\documentclass[12pt]{article}
\usepackage[utf8]{inputenc}
\usepackage[T1]{fontenc}
\usepackage[margin=1in]{geometry}
\usepackage{setspace}
\usepackage{xcolor}
\usepackage[colorlinks=true, urlcolor=blue!60!black, linkcolor=blue!60!black]{hyperref}
\setlength{\parindent}{0em}
\setlength{\parskip}{0.45em}
\renewcommand{\baselinestretch}{1.0}
\begin{document}
\begin{center}
{\Large \textit{FDA v. Alliance for Hippocratic Medicine}, 602 U.S. 367 (2024)\par}
\vspace{0.3em}
{\normalsize Case Brief --- One-Page Edition\par}
\end{center}

\textbf{Memory Jogger.} Pro-life doctors who neither prescribe nor use mifepristone could not challenge its deregulation; there is no ``doctor standing'' (\href{https://www.supremecourt.gov/opinions/23pdf/602us1r35_bqmd.pdf}{602 U.S. 367}).

\textbf{Facts.} In 2016 and 2021 the FDA eased mifepristone's restrictions, extending use to ten weeks, allowing non-physician prescribers, and permitting mail dispensing. Pro-life medical associations and doctors sued under the APA, although none prescribed or used the drug and the FDA regulated nothing they did.

\textbf{Procedural History.} The district court effectively pulled the drug; the Supreme Court stayed that order; the Fifth Circuit upheld relief against the 2016 and 2021 actions.

\textbf{Votes.} Reversed, 9--0: Kavanaugh for the Court; Thomas concurred.

\textbf{Holding.} Unregulated plaintiffs lack standing to challenge the deregulation of a drug they do not prescribe or use.

\textbf{Reasoning.} Courts are not a forum for bystanders' policy objections (``What's it to you?''), and causation is hard to show when the government regulates someone else. Federal conscience laws protect doctors from participating, breaking the causal chain; downstream costs of treating complications are speculative and attenuated, and doctor standing would have no limit; and \textit{Havens Realty} reaches only direct interference with an organization's core activities, not advocacy spending. That no one else may have standing is no reason to find it. Thomas urged reconsidering associational standing.

\textbf{Examples.} \emph{Same-Side:} advocacy groups fighting a voter-roll law. \emph{Opposite-Side:} a rival drug maker. \emph{Fence-Sitter:} a legal-aid group barred from detained clients.

\textbf{Critique.} Praised as disciplined, it left the merits open to state suits after \textit{Loper Bright}; critics contrast looser standing in \textit{Biden v. Nebraska} and warn it may hobble civil-rights nonprofits.

\textbf{Key Quotation.} ``But an organization that has not suffered a concrete injury caused by a defendant's action cannot spend its way into standing simply by expending money to gather information and advocate against the defendant's action.'' --- Justice Kavanaugh, 602 U.S. at 394.

\end{document}
